\documentclass[../DoAn.tex]{subfiles}
\begin{document}

\section*{Quy định chung}
Dưới đây là một số quy định và hướng dẫn viết đồ án tốt nghiệp mà bắt buộc sinh viên phải đọc kỹ và tuân thủ nghiêm ngặt.

Sinh viên cần đảm bảo tính thống nhất toàn báo cáo (font chữ, căn dòng hai bên, hình ảnh, bảng, margin trang, đánh số trang, v.v.). Để làm được như vậy, sinh viên chỉ cần sử dụng các định dạng theo đúng template ĐATN này. Khi paste nội dung văn bản từ tài liệu khác của mình, sinh viên cần chọn kiểu Copy là “Text Only” để định dạng văn bản của template không bị phá vỡ/vi phạm.

Tuyệt đối cấm sinh viên đạo văn. Sinh viên cần ghi rõ nguồn cho tất cả những gì không tự mình viết/vẽ lên, bao gồm các câu trích dẫn, các hình ảnh, bảng biểu, v.v. Khi bị phát hiện, sinh viên sẽ không được phép bảo vệ ĐATN.

Tất cả các hình vẽ, bảng biểu, công thức, và tài liệu tham khảo trong ĐATN nhất thiết phải được SV giải thích và tham chiếu tới ít nhất một lần. Không chấp nhận các trường hợp sinh viên đưa ra hình ảnh, bảng biểu tùy hứng và không có lời mô tả/giải thích nào.

Sinh viên tuyệt đối không trình bày ĐATN theo kiểu viết ý hoặc gạch đầu dòng. ĐATN không phải là một slide thuyết trình; khi người đọc không hiểu sẽ không có ai giải thích hộ. Sinh viên cần viết thành các đoạn văn và phân tích, diễn giải đầy đủ, rõ ràng. Câu văn cần đúng ngữ pháp, đầy đủ chủ ngữ, vị ngữ và các thành phần câu.

Khi thực sự cần liệt kê, sinh viên nên liệt kê theo phong cách khoa học với các ký tự La Mã. Ví dụ, nhiều sinh viên luôn cảm thấy hối hận vì (i) chưa cố gắng hết mình, (ii) chưa sắp xếp thời gian học/chơi một cách hợp lý, (iii) chưa tìm được người yêu để chia sẻ quãng đời sinh viên vất vả, và (iv) viết ĐATN một cách cẩu thả.

Trong một số trường hợp nhất thiết phải dùng các nhóm để phân tách nội dung, sinh viên cần thống nhất quy cách thể hiện cho toàn bộ báo cáo, hạn chế tối đa việc lạm dụng các ký tự chấm tròn tự do, thay vào đó nên diễn dịch thành văn bản liên tục nhằm giữ vững tính học thuật và tâm huyết trong từng trang báo cáo ĐATN.

\section{Ngành học}
Sinh viên lưu ý viết đúng ngành hoặc chuyên ngành trên bìa và trên gáy theo đúng quy định hiện hành của Trường Công nghệ Thông tin và Truyền thông. Tên ngành đào tạo chính thức phụ thuộc vào chương trình học tập mà sinh viên đăng ký và được lưu trữ trên hệ thống quản lý dữ liệu. Sinh viên có thể truy cập tài khoản cá nhân để tra cứu chính xác thông tin. 

Nhằm đảm bảo tính chuẩn hóa, một số mẫu tên ngành đào tạo chuẩn được phân định rõ ràng bao gồm: (i) đối với hệ kỹ sư chính quy từ khóa K61 trở về trước sử dụng tên ngành Kỹ thuật phần mềm, trong khi từ khóa K62 trở về sau đổi tên thành ngành Khoa học máy tính; (ii) đối với hệ cử nhân công nghệ thông tin sử dụng tên gọi ngành Công nghệ thông tin; (iii) đối với các chương trình đào tạo tiên tiến EliteTech bao gồm hệ Việt Nhật và Kỹ sư tài năng (KSTN) sử dụng tên ngành Công nghệ thông tin, chương trình ICT Global sử dụng tên tiếng Anh Information Technology, và chương trình DS\&AI áp dụng tên ngành Khoa học dữ liệu.

\section{Đánh dấu và đánh số danh sách}
Việc sử dụng danh sách trong môi trường LaTeX khá đơn giản và không yêu cầu sinh viên phải tích hợp thêm bất kỳ gói bổ sung nào. Hệ thống cung cấp hai môi trường định dạng danh mục chính bao gồm:

Môi trường thứ nhất là kiểu liệt kê không có thứ tự (\texttt{itemize}). Để triển khai cấu trúc danh mục đánh dấu này, mã nguồn được khai báo tuần tự theo cú pháp:
\begin{verbatim}
\begin{itemize}
    \item Nội dung thứ nhất được viết ở đây.
    \item Nội dung thứ hai được viết ở đây.
\end{itemize}
\end{verbatim}

Môi trường thứ hai là kiểu liệt kê có thứ tự theo cấp bậc (\texttt{enumerate}). Để sử dụng kiểu định dạng đánh số hệ thống, cú pháp mã nguồn được thiết lập như sau:
\begin{verbatim}
\begin{enumerate}
    \item Nội dung thứ nhất được viết ở đây.
    \item Nội dung thứ hai được viết ở đây.
\end{enumerate}
\end{verbatim}

Toàn bộ các nội dung trình bày trong cả hai môi trường liệt kê bắt buộc phải đặt ngay sau lệnh \texttt{\textbackslash item}. Ngoài ra, LaTeX còn hỗ trợ một số kiểu tùy biến danh mục nâng cao khác, sinh viên có thể chủ động tham khảo chi tiết tại tài liệu hướng dẫn trực tuyến của Overleaf.

\section{Cách thêm bảng biểu}
\begin{table}[ht!]
\centering
\small
\renewcommand{\arraystretch}{1.2}
\begin{tabular}{||c c c c||} 
 \hline
 Col1 & Col2 & Col2 & Col3 \\ [0.5ex] 
 \hline\hline
 1 & 6 & 87837 & 787 \\ 
 2 & 7 & 78 & 5415 \\
 3 & 545 & 778 & 7507 \\
 4 & 545 & 18744 & 7560 \\
 5 & 88 & 788 & 6344 \\ [1ex] 
 \hline
\end{tabular}
\caption{Bảng mẫu thực nghiệm kiểm thử tham số và nhãn định danh}
\label{table:example_appendix}
\end{table}

Bảng \ref{table:example_appendix} là một ví dụ trực quan về quy cách xây dựng và định dạng bảng dữ liệu trong báo cáo. Theo quy định, tất cả các bảng biểu xuất hiện trong đồ án tốt nghiệp nhất thiết phải được dẫn dắt, phân tích và bình luận nội dung chi tiết trong phần văn bản liên quan. Do việc thiết lập cấu trúc bảng trong môi trường LaTeX tương đối phức tạp, sinh viên được khuyến khích sử dụng các công cụ hỗ trợ sinh mã tự động trực tuyến để tối ưu hóa thời gian xây dựng mã nguồn.

\section{Chèn hình ảnh minh họa}
\begin{figure}[ht!]
\centering
\includegraphics[width=0.65\linewidth]{Hinhve/IoT.png}
\caption{Mô hình kết nối kiến trúc Internet vạn vật (IoT)}
\label{fig:iot_appendix}
\end{figure}

Hình \ref{fig:iot_appendix} mô tả cấu trúc ví dụ về quy trình chèn hình ảnh kỹ thuật số vào tài liệu. Phần chú thích giải nghĩa của hình vẽ (\texttt{caption}) bắt buộc phải được đặt cố định ở vị trí ngay dưới hình ảnh tương ứng. Sinh viên lưu ý toàn bộ hệ thống hình vẽ được sử dụng trong báo cáo đều phải có nhãn định danh (\texttt{label}), phải được tham chiếu trực tiếp ít nhất một lần trong phần nội dung luận văn và đi kèm các đoạn văn phân tích, diễn giải khoa học cụ thể.

\section{Tài liệu tham khảo và trích dẫn}
\paragraph{Cách thức liệt kê trích dẫn}\mbox{}

Hệ thống báo cáo đồ án tốt nghiệp áp dụng đồng bộ quy cách liệt kê tài liệu tham khảo và trích dẫn theo tiêu chuẩn quốc tế IEEE. Minh họa cho việc thực hiện trích dẫn mã nguồn trong đoạn văn được thể hiện cụ thể thông qua lệnh gọi \cite{scott2013sdn}, hoặc cấu hình tham chiếu phối hợp như \cite{ashton2009internet}. Hệ thống biên dịch sẽ tự động lọc và chỉ hiển thị trong danh mục tài liệu tham khảo cuối quyển những nguồn dữ liệu thực sự có sử dụng lệnh trích dẫn \texttt{\textbackslash cite\{\}} trong nội dung văn bản. Toàn bộ tài liệu tham khảo phải có nguồn gốc phát hành rõ ràng, đáng tin cậy và hạn chế tối đa việc sử dụng các trang thông tin cộng đồng hoặc bách khoa toàn thư mở.

\paragraph{Phân loại các nguồn tài liệu}\mbox{}

Các nguồn tài liệu tham khảo hợp quy chuẩn khoa học được phân định thành các nhóm chính bao gồm: hệ thống sách giáo trình chuyên ngành, các bài báo khoa học xuất bản trên các tạp chí uy tín, báo cáo chuyên môn tại các hội nghị khoa học quốc tế và các tài liệu đặc tả kỹ thuật chính thống được công bố trên Internet bởi các tổ chức công nghệ toàn cầu.

\section{Phương trình và công thức toán học}
Hạ tầng mã nguồn của hệ thống đã tích hợp sẵn các gói thư viện toán học lõi bao gồm \texttt{amsmath}, \texttt{amssymb} và \texttt{amsfonts}. Quy trình khởi tạo một phương trình toán học có gán nhãn định danh tự động được thể hiện thông qua cấu trúc Phương trình (\ref{eq:integral_demo}):

\begin{equation}\label{eq:integral_demo}
    F(x) = \int^a_b \frac{1}{3}x^3
\end{equation}

Phương trình (\ref{eq:integral_demo}) đại diện cho mô hình tính toán tích phân cơ bản. Trong trường hợp toán tử hiển thị là các công thức bổ trợ thuần túy và không cần thiết phải đánh số thứ tự quản lý, sinh viên áp dụng cấu hình toán ký tự tự do như sau:

\begin{equation*}
    x[t_n] = \frac{1}{\sqrt{N}} \sum_{k=0}^{N-1}X[f_k]e^{j 2\pi n k/N}
\end{equation*}

Công thức toán học trên biểu diễn cấu trúc toán tử của phép biến đổi Fourier rời rạc ngược (IDFT) trong phân tích xử lý tín hiệu.

\section{Quy cách đóng quyển báo cáo}
Phần bìa trước của đồ án tốt nghiệp phải được chế bản và in ấn chính xác theo đúng khuôn mẫu quy định của nhà trường. Toàn bộ phần bìa trước và bìa sau bắt buộc phải sử dụng chất liệu giấy liền khổ, thực hiện quy trình dán gáy bằng công nghệ đóng keo nhiệt chuyên dụng, tuyệt đối cấm sử dụng các hình thức cố định thô sơ bằng băng dính đóng sách hoặc dập ghim vật lý thông thường.

Phần gáy của quyển đồ án tốt nghiệp yêu cầu phải in rõ ràng và đầy đủ các thông tin tóm tắt bao gồm: Kỳ học hoàn thành đồ án, Tên ngành đào tạo chính thức, Họ và tên khai sinh của sinh viên, và Mã số sinh viên duy nhất. Cấu trúc nội dung hiển thị trên gáy sách được quy định mẫu chuẩn: \texttt{2022.1 - KỸ THUẬT MÁY TÍNH - NGUYỄN VĂN A - 20221234}. Sinh viên cần rà soát kỹ lưỡng để đảm bảo toàn bộ gáy quyển tuân thủ đúng định dạng khoảng cách và cỡ chữ quy chuẩn trước khi nộp bản cứng về văn phòng Viện.

\begin{figure}[ht!]
\centering
\includegraphics[width=0.5\linewidth]{Hinhve/Bia.PNG}
\caption{Quy cách kỹ thuật đóng quyển báo cáo đồ án tốt nghiệp chuẩn}
\label{fig:Bia}
\end{figure} 

\begin{figure}[ht!]
\centering
\includegraphics[width=0.5\linewidth]{Hinhve/GayBia.PNG}
\caption{Quy cách định dạng cấu trúc font chữ in trên gáy bìa cứng}
\label{fig:GayBia} % Đã sửa thành fig:GayBia để triệt tiêu hoàn toàn lỗi trùng nhãn
\end{figure} 

\end{document}